\section{Related Work}\label{sec:related}
\subsection{Conversational tools and smart-home data}
ReAct demonstrates how a language model can obtain information by interleaving reasoning with external actions~\cite{react}. AdaHome routes smart-home requests and uses interaction history to handle preference-dependent commands~\cite{adahome}. DietSense adopts the separation between interpreting a request and executing a tool. A preference about drying clothes indoors may guide a recommendation, whereas a statement that the dehumidifier ran requires an operation record or confirmation. Keeping these histories separate allows a resident to correct one without changing the other.

AgentSense generates labeled sensor streams from simulated household activities~\cite{agentsense}. Such streams can support activity-recognition research. In the present setting, the difficult cases are episodes in which several devices operate together or a resident cannot recall what happened. A simulation supplies its own event labels. A home deployment must account for who confirmed an event and when that information became available.

\subsection{Environmental models and the choice of labels}
The CBE Thermal Comfort Tool implements comfort models with explicit environmental and personal inputs~\cite{cbe}. DietSense has a single temperature and humidity sensor. It can report a trend and ask how it felt, but lacks several inputs needed for a full comfort assessment. Likewise, the office dataset studied by Candanedo and Feldheim associates environmental measurements with occupancy~\cite{occupancy}. A successful occupancy prediction would not establish which device was operating. Table~\ref{tab:work} identifies the evidence each source contributes and the question left for the pilot.

\begin{table*}[t]
\caption{Prior work and its applicability to the single-room pilot}
\label{tab:work}\centering\footnotesize
\setlength{\tabcolsep}{3pt}
\begin{tabularx}{\textwidth}{@{}l l Y Y Y Y Y Y@{}}
\toprule
Ref. & Year & Problem & Method & Domain knowledge & Data or input & Contribution & Limit for this task\\
\midrule
\cite{react} & 2023 & Retrieve evidence during a task & Tool-using language model & Interactive reasoning & Question answering and decision tasks & External results inform subsequent actions & Sensor location and time need explicit checks\\
\cite{adahome} & 2026 & Interpret household requests & Local language model and routing & Home preferences & Commands and preference histories & Routes requests by intent & A preference cannot confirm a device state\\
\cite{agentsense} & 2026 & Obtain activity labels & Agent-based simulation & Home activities & Virtual sensor streams & Produces labeled activity sequences & Simulated labels cannot verify a real event\\
\cite{deDear} & 1998 & Explain variation in comfort & Adaptive model & Thermal preference & Field surveys in buildings & Relates context to preferred conditions & A population trend cannot label one episode\\
\cite{occupancy} & 2016 & Estimate occupancy & Statistical learning & Office environment & Sensors and occupancy labels & Predicts whether an office is occupied & Occupancy and appliance use are distinct targets\\
\cite{cbe} & 2020 & Calculate comfort indices & Physical and empirical models & Thermal comfort & Environmental and personal inputs & Makes model assumptions explicit & Some inputs are unmeasured in the pilot\\
\bottomrule
\end{tabularx}
\end{table*}

